% !PRINTED 65
\weberpage{65}
But by (10) we have
\[
2a_2n_1n_2\omega_1 = a_1n_1n_2\omega_1 + a_2n_1n_2\omega_2 + mn_1n_2,
\]
and accordingly this last formula becomes:
\weq{16}{\begin{aligned}
&T(u+n_1\omega_1+n_2\omega_2)\\
&= e^{-\pi i[(n_1a_1+n_2a_2)(2u+n_1\omega_1+n_2\omega_2)+\pi i(b_1n_1+b_2n_2+mn_1n_2)]}T(u).
\end{aligned}}
% [sic] German source includes a nested \pi i factor in the final bracket.
And here $n_1$ and $n_2$ are arbitrary positive or negative integers.

If two $T$-functions of orders $m$ and $m'$, with the characters $\gamma$, $\gamma'$, are multiplied together, there arises a new $T$-function whose order is $m+m'$ and whose character is $\gamma+\gamma'$. If $(g_1, g_2)$ and $(g_1', g_2')$ are the characteristics of the two factors, then $(g_1+g_1', g_2+g_2')$ is the characteristic of the product.

\begin{center}
\textbf{\S\,18. \gesperrt{Relations between related $T$-Functions.}}
\end{center}

Two $T$-functions with the same periods, the same order and the same character we shall call \gesperrt{related}. If $T'(u)$ is a $T$-function related to $T(u)$, and if $a_1'$, $a_2'$, $b_1'$, $b_2'$ have the same meaning for $T'$ as $a_1$, $a_2$, $b_1$, $b_2$ have for $T$, then in consequence of the assumed relationship [§\,17, (10), (11)]:
\weq{1}{a_1'\omega_2 - a_2'\omega_1 = a_1\omega_2 - a_2\omega_1,}
\weq{2}{b_1'\omega_2 - b_2'\omega_1 = b_1\omega_2 - b_2\omega_1 + 2n_1\omega_2 - 2n_2\omega_1,}
where $n_1$, $n_2$ are integers.

Accordingly $\lambda$ and $\mu$ can be determined by § 17, (9) in such a way that the periodicity factors of the function
\[
e^{-\pi i(\lambda u^2+\mu u)}T(u)
\]
become the same as those of the function $T'(u)$, and hence we have the theorem:

1. \gesperrt{Related $T$-functions can, by the addition of a $T$-function of order zero:}
\[
L = Ce^{-\pi i(\lambda u^2+\mu u)}
\]
\gesperrt{as a factor, be given the same periodicity factors.}

Further, from the equality of the characters there follows immediately:

2. \gesperrt{If related $T$-functions of order $m$ have $m-1$ common zero points in the period parallelogram, then they have the $m$th in common as well;} and from this:
%CONT>>
